% ****************************************************************

\subsection{SplitPorts and SplitVector}

\label{sec-SplitPorts-package}

\index{SplitPorts@\te{SplitPorts}}
\index{SplitPorts@\te{SplitVector}}

NOTE: \hfill \fbox{\begin{minipage}{0.9\textwidth}\footnotesize
    \te{SplitPorts} is a new feature, introduced in release 2026.07.  We
    welcome feedback on its usability and on this documentation.
\end{minipage}}

{\footnotesize Although aimed at Verilog, BlueSim also treats split
  ports in the same way. In particular, it will generate waveform
  files with the same ports/buses seen in Verilog.}

% ----------------------------------------------------------------

\subsubsection{Description: Purpose of ``port splitting''}

\index{Port splitting}
\index{Port splitting|seealso{\te{SplitPorts}}}

{\footnotesize (In this section we will use the terms ``port'' and ``bus'' interchangeably.)}

Consider a {\BHb} or {\BSVb} module compiled by \emph{bsc} to
Verilog. Suppose an interface method has a ``struct'' argument or
result, becoming a Verilog input or output port, respectively.  These
appear in Verilog as flat bit-vectors (usually the concatenation of
the bit representations of the fields).  These are difficult to use in
several RTL-level situations:

\begin{itemize}
  \item Viewing waveforms (whether generated from Verilog simulation or from Bluesim).
  \item Interfacing with existing Verilog code.
\end{itemize}

In both cases, we typically want a separate, named bus for each field
of the struct, instead of having to visualize relevant bits of the
full flat bit-vector.  \te{SplitPorts} provides this capability, and
more:

\begin{itemize}

  \item Splitting can be recursive, on nested types: struct, vectors
    and tuples, and algebraic types (tagged unions).

  \item Defaults make it very easy to split a structured type
    ``canonically'' (into its fields). To split in more custom ways
    (\emph{e.g.,} combining splits, sub-splitting fields,
    \emph{etc.}), it is easy to override defaults.

    NOTE: \hfill \fbox{\begin{minipage}{0.8\textwidth}\footnotesize
      We do not enable splitting by default on the polymorphic library
      types Vectors and Tuples.
    \end{minipage}}

\end{itemize}

All this is achieved with library type classes---no new syntax or
annotations--- and is equally available in {\BSVb} and {\BHb}.

% ----------------

\paragraph{Canonical \emph{vs.} Custom splitting}

\hspace*{1em}

\index{SplitPorts@\te{SplitPorts}!Canonical port splitting}
\index{Canonical port splitting}
\index{Canonical port splitting|seealso{\te{SplitPorts}}}

We use the term ``canonical splitting'' where each struct field
becomes a port, one-to-one, each vector/tuple element becomes a port,
one-to-one, and so on.  For all other kinds of splitting we use the
term ``custom splitting''.

Splitting of a usertype $t$ needs to be specified by the user.
Defaults and utilities simplify this for canonical splitting.

We expect that canonical splitting is what you'd want most of the
time.  Custom splitting can also be useful in situations like these:

\begin{itemize}

  \item The desired Verilog ports are not one-to-one with the
    {\BLANGS} structured type's fields.

  \item To ``impedance-match'' {\BLANGS} coding to some other coding
    in external Verilog or other IP, such as ``little-endian'' to
    ``big-endian'', or SystemVerilog's ``packed'' structures.

\end{itemize}

% ----------------

\paragraph{Shallow \emph{vs.} Deep Splitting}

\hspace*{1em}

\index{SplitPorts@\te{SplitPorts}!Shallow port splitting}
\index{Shallow port splitting}
\index{Shallow port splitting|seealso{\te{SplitPorts}}}

\index{SplitPorts@\te{SplitPorts}!Deep port splitting}
\index{Deep port splitting}
\index{Deep port splitting|seealso{\te{SplitPorts}}}

In general, a type can have \emph{nested} structure.  For example, in
a struct type $t_1$ a field may have struct type $t_2$, or a vector
type, or a tuple type, or an algebraic type (tagged union).  And,
recursively, those types may in turn contain further nested types.

When an interface-method arg or result has such a type, we can split
it to various levels.

A \emph{shallow split} only splits the top-level type of the nested
structure, \emph{e.g.,} a top-level sruct into ports for its fields.
The fields themselves are not split, their ports remain as aggregate
bit-vectors.

A \emph{deep split}, on the other hand, splits the fields, and fields
of fields, and so on, recursively, down to ``leaf-level'' fields
(scalars, which do not themselves have multiple components).

% ----------------

\paragraph{Examples}

\hspace*{1em}

Several examples are given below, in
Sections~\ref{sec:splitports_examples_BH}
and
Sections~\ref{sec:splitports_examples_BSV}.

% ----------------------------------------------------------------

\subsubsection{Packages} 

\label{sec:Library_SplitPorts}

The {\tt SplitPorts} type class and {\tt Port} type are defined in the
Standard Prelude which is always implicitly imported.  They are
described in Sections~\ref{sec-SplitPorts-Port-type} and
\ref{sec-SplitPorts-typeclass}, respectively.  They are adequate for
``manually'' specifying splitting.

The {\tt SplitPorts} and {\tt SplitVector} packages must be imported
for additional facilities (simplified canonical splitting and
splitting of vectors). Detailed usage examples are shown in sections
below.

\begin{center}
  \index{SplitPorts@\te{SplitPorts}!\te{SplitPorts} package}
  \index{SplitVector@\te{SplitVector}!\te{SplitVector} package}
  \begin{minipage}[t]{0.3\textwidth}
    {\BHb} syntax:

    \begin{tabular}{|l|}
      \hline
      {\tt import SplitPorts} \\
      \hline
      {\tt import SplitVector} \\
      \hline
    \end{tabular}
  \end{minipage}
  \hspace*{2em}
  \begin{minipage}[t]{0.3\textwidth}
    {\BSVb} syntax:

    \begin{tabular}{|l|}
      \hline
      {\tt import SplitPorts :: *;} \\
      \hline
      {\tt import SplitVector :: *;} \\
      \hline
    \end{tabular}
  \end{minipage}
\end{center}

% ----------------------------------------------------------------

\subsubsection{Typeclasses}

\label{sec:splitports_shallow_deep}

\index{ShallowSplitPorts@\te{ShallowSplitPorts} (package for \te{SplitPorts})}
\index{SplitPorts@\te{SplitPorts}!Related type class \te{ShallowSplitPorts}}
\index[typeclass]{ShallowSplitPorts}

\index{shallowSplitPorts@\te{shallowSplitPorts}
       (function in \te{ShallowSplitPorts} type class}
\index{shallowUnsplitPorts@\te{shallowUnsplitPorts}
       (function in \te{ShallowSplitPorts} type class}
\index{shallowSplitPortNames@\te{shallowSplitPortNames}
       (function in \te{ShallowSplitPorts} type class}


\index{DeepSplitPorts@\te{DeepSplitPorts} (package for \te{SplitPorts})}
\index{SplitPorts@\te{SplitPorts}!Related type class \te{DeepSplitPorts}}
\index[typeclass]{DeepSplitPorts}

\index{deepSplitPorts@\te{deepSplitPorts}
       (function in \te{DeepSplitPorts} type class}
\index{deepUnsplitPorts@\te{deepUnsplitPorts}
       (function in \te{DeepSplitPorts} type class}
\index{deepSplitPortNames@\te{deepSplitPortNames}
       (function in \te{DeepSplitPorts} type class}


The following type classes are defined in the {\tt SplitPorts} library
package. The type class methods simplify the specification of
canonical splitting (see examples in
Section~\ref{sec:splitports_examples_BH_simplified}).

\begin{quote}
\begin{verbatim}
class ShallowSplitPorts a p | a -> p where
  type ShallowPortsOf a = p
  shallowSplitPorts :: a -> p
  shallowUnsplitPorts :: p -> a
  shallowSplitPortNames :: a -> String -> List String
\end{verbatim}
\end{quote}

\begin{quote}
\begin{verbatim}
class DeepSplitPorts a p | a -> p where
  type DeepPortsOf a = p
  deepSplitPorts :: a -> p
  deepUnsplitPorts :: p -> a
  deepSplitPortNames :: a -> String -> List String
\end{verbatim}
\end{quote}

NOTE: \hfill \fbox{\begin{minipage}{0.9\textwidth}\footnotesize

  In the library these definitions are in {\BHb} syntax but they can
  be imported and used equally in {\BSVb} code.

\end{minipage}}

%  ----------------------------------------------------------------

\subsubsection{Examples in {\BHb} of splitting user-defined types}

\label{sec:splitports_examples_BH}

{\footnotesize
For {\BSVb} versions of these, please see Section~\ref{sec:splitports_examples_BSV}.}

% ----------------

\paragraph{Example ({\BHb}): Baseline (before splitting) for port splitting examples}

\label{sec:splitports_examples_BH_baseline}

\hspace*{1em}

Consider the following struct and interface types:
\begin{quote}
\begin{verbatim}
struct Foo =
  x :: Int 8
  y :: Int 8
 deriving (Bits)    -- Represented as Bit 16

struct Bar =
  v :: Vector 3 Bool
  w :: (Bool, UInt 16)
  z :: Foo
 deriving (Bits)    -- Represented as Bit 36    (3+1+16+16)

interface SplitTest =
  mv  :: Foo -> Bar
  mav :: Foo -> ActionValue Bar
\end{verbatim}
\end{quote}

If we instantiate a module {\tt mod} with this interface, the methods
may be invoked like this:

\begin{tabbing}
\hspace*{4em} \= \hspace*{2em} \= \kill
\hspace*{4em} \>               \> {\tt let (b1::Bar)  = mod.mv  (a1::Foo)} \\
              \>  or           \> {\tt     (b2::Bar) <- mod.mav (a2::Foo)}
\end{tabbing}

In the RTL, the module will have the following ports:

\begin{center}
  \begin{tabular}{|c||c|}
    \hline
     Output &
     Input  \\
    \hline
     {\tt mv}    &
     {\tt mv\_1} \\
     36 bits &
     16 bits \\
    \hline
     {\tt mav}    &
     {\tt mav\_1} \\
     36 bits &
     16 bits \\
    \hline
\end{tabular}

{\footnotesize (The ``{\tt \_1}'' suffix on the args are for ``first argument''.)}
\end{center}

%  ----------------------------------------------------------------

\paragraph{Example (\BH): shallow-split {\tt Foo}}

\label{sec:splitports_examples_BH_Foo_Shallow}

\hspace*{1em}

{\footnotesize (This example does not need import of {\tt SplitPorts} package.)}

Let us add the following declaration to our program, making {\tt Foo}
an instance of the {\tt SplitPorts} type class:

\begin{quote}
\begin{verbatim}
type T2 = (Port (Int 8), Port (Int 8))

instance SplitPorts Foo T2 where
  splitPorts f = (Port (f.x), Port (f.y))
  unsplitPorts (Port x, Port y) = Foo {x=x; y=y}
  portNames _ base = Cons (base +++ "_x")
                       (Cons (base +++ "_y") Nil)
\end{verbatim}
\end{quote}

{\tt splitPorts} takes an argument {\tt f} of type {\tt Foo}, and
returns a 2-tuple of port descriptors, one per field of {\tt Foo}.
Each component applies the constructor {\tt Port} to the corresponding
component.

{\tt unsplitPorts} is the inverse: it takes a 2-tuple of port
descriptors and recombines them into a {\tt Foo} value.

{\tt portNames} returns a list of desired names for the ports, by
appending suffixes to bhe base name.

In the RTL, the module will have the following ports:

\begin{center}
  \begin{tabular}{|c||c|c|}
    \hline
     Output &
     Input  &
     Input  \\
    \hline
     {\tt mv}       &
     {\tt mv\_1\_x} &
     {\tt mv\_1\_y} \\
     36 bits &
      8 bits &
      8 bits \\
    \hline
     {\tt mav}       &
     {\tt mav\_1\_x} &
     {\tt mav\_1\_y} \\
     36 bits &
      8 bits &
      8 bits \\
    \hline
  \end{tabular}
\end{center}

The argument ports, of type {\tt Foo} have been split.

%  ----------------------------------------------------------------

\paragraph{Example (\BH): shallow-split {\tt Foo}, shallow-split {\tt Bar}}

\label{sec:splitports_examples_BH_Foo_Shallow_Bar_Shallow}

\hspace*{1em}

{\footnotesize (This example does not need import of {\tt SplitPorts} package.)}

Suppose we add the following declaration to our program, making {\tt
  Bar} also an instance of the {\tt SplitPorts} type class:

\begin{quote}
\begin{verbatim}
type T6 = (Port Bool, Port Bool, Port Bool,
           Port Bool, Port (UInt 16),
           Port Foo)

instance SplitPorts Bar T6 where
  splitPorts b = let (w1,w2) = b.w
                 in (Port (b.v !! 0), Port (b.v !! 1), Port (b.v !! 2),
                     Port (w1), Port (w2),
                     Port (b.z))

  unsplitPorts (Port v0, Port v1, Port v2, Port w1, Port w2, Port z) =
    Bar {v = v0 :> v1 :> v2 :> nil;
         w = (w1, w2);
         z = z}

  portNames _ base = Cons (base +++ "_v_0")
                       (Cons (base +++ "_v_1")
                          (Cons (base +++ "_v_2")
                             (Cons (base +++ "_w_1")
                                (Cons (base +++ "_w_2")
                                   (Cons (base +++ "_z")
                                      Nil)))))
\end{verbatim}
\end{quote}

In the RTL, the module will now have the following ports:

\begin{center}
  \begin{tabular}{|c|c|c|c|c|c||c|c|}
    \hline
     Output &
     Output &
     Output &
     Output &
     Output &
     Output &
     Input  &
     Input  \\
    \hline
     {\tt mv\_v\_0}  &
     {\tt mv\_v\_1}  &
     {\tt mv\_v\_2}  &
     {\tt mv\_w\_1}  &
     {\tt mv\_w\_2}  &
     {\tt mv\_z}     &
     {\tt mv\_1\_x}  &
     {\tt mv\_1\_y} \\
      1 bit  &
      1 bit  &
      1 bit  &
      1 bit  &
     16 bits &
     16 bits &
      8 bits &
      8 bits \\
    \hline
     {\tt mav\_v\_0}  &
     {\tt mav\_v\_1}  &
     {\tt mav\_v\_2}  &
     {\tt mav\_w\_1}  &
     {\tt mav\_w\_2}  &
     {\tt mav\_z}     &
     {\tt mav\_1\_x}  &
     {\tt mav\_1\_y} \\
      1 bit  &
      1 bit  &
      1 bit  &
      1 bit  &
     16 bits &
     16 bits &
      8 bits &
      8 bits \\
    \hline
  \end{tabular}
\end{center}

Even though each output port, of type {\tt Bar}, contains a field
``{\tt z}'' of type {\tt Foo}, that nested field is not split.
Splitting is only applied to the top-level type ({\tt Bar}), and in
its {\tt splitPorts} function we chose not to split the nested field
{\tt z}.

%  ----------------------------------------------------------------

\paragraph{Example (\BH): shallow-split {\tt Foo}, deep-split {\tt Bar}}

\label{sec:splitports_examples_BH_Foo_Shallow_Bar_Deep}

\hspace*{1em}

{\footnotesize (This example does not need import of {\tt SplitPorts} package.)}

Suppose we change the declaration of the {\tt Bar} instance of {\tt
SplitPorts} to the following:

\begin{quote}
\begin{verbatim}
type T7 = (Port Bool, Port Bool, Port Bool,
           Port Bool, Port (UInt 16),
           Port (Int 8), Port (Int 8))


instance SplitPorts Bar T7 where
  splitPorts b = let (w1,w2) = b.w
                 in (Port (b.v !! 0), Port (b.v !! 1), Port (b.v !! 2),
                     Port (w1), Port (w2),
                     Port (b.z.x), Port (b.z.y))

  unsplitPorts (Port v_0, Port v_1, Port v_2,
                
                Port w_1, Port w_2,
                
                Port z_x, Port z_y)
   = Bar {v = v_0 :> v_1 :> v_2 :> nil;
          w = (w_1, w_2);
          z = Foo {x = z_x; y = z_y}}

  portNames _ base = Cons (base +++ "_v_0")
                       (Cons (base +++ "_v_1")
                          (Cons (base +++ "_v_2")
                             (Cons (base +++ "_w_1")
                                (Cons (base +++ "_w_2")
                                   (Cons (base +++ "_z_x")
                                      (Cons (base +++ "_z_y")
                                         Nil))))))
\end{verbatim}
\end{quote}

In the RTL, the module will now have the following ports:

\begin{tabular}{|c|c|c|c|c|c|c||c|c|}
    \hline
     Output &
     Output &
     Output &
     Output &
     Output &
     Output &
     Output &
     Input  &
     Input  \\
    \hline
     {\tt mv\_v\_0}  &
     {\tt mv\_v\_1}  &
     {\tt mv\_v\_2}  &
     {\tt mv\_w\_1}  &
     {\tt mv\_w\_2}  &
     {\tt mv\_z\_x}  &
     {\tt mv\_z\_y}  &
     {\tt mv\_1\_x}  &
     {\tt mv\_1\_y} \\
      1 bit  &
      1 bit  &
      1 bit  &
      1 bit  &
     16 bits &
      8 bits &
      8 bits &
      8 bits &
      8 bits \\
    \hline
     {\tt mav\_v\_0}  &
     {\tt mav\_v\_1}  &
     {\tt mav\_v\_2}  &
     {\tt mav\_w\_1}  &
     {\tt mav\_w\_2}  &
     {\tt mav\_z\_x}  &
     {\tt mav\_z\_y}  &
     {\tt mav\_1\_x}  &
     {\tt mav\_1\_y} \\
      1 bit  &
      1 bit  &
      1 bit  &
      1 bit  &
     16 bits &
      8 bits &
      8 bits &
      8 bits &
      8 bits \\
    \hline
\end{tabular}

Now note that the nested field {\tt z} of type {\tt Bar} on the output
is also split, because we specfied that in the {\tt splitPorts}
function.

%  ----------------------------------------------------------------

\paragraph{Example (\BH): custom-split {\tt Foo}}

\label{sec:splitports_examples_BH_Custom}

\hspace*{1em}

{\footnotesize (This example does not need import of {\tt SplitPorts} package.)}

(Compare this with shallow-split {\tt Foo} in
Section:~\ref{sec:splitports_examples_BH_Foo_Shallow}.)

Suppose we change the declaration of the {\tt Foo} instance of {\tt
SplitPorts} to the following:

\begin{quote}
\begin{verbatim}
type T3 = (Port (Int 8), Port Bool, Port (Bit 7))

instance SplitPorts Foo T3 where

  splitPorts f = (Port f.x,
                  Port (f.y > 0), Port $ truncate $ pack f.y)

  unsplitPorts (Port x, Port s, Port y) =
      Foo { x = x;
            y = (if s then id else negate) $ unpack $ zeroExtend y; }

  portNames _ base = Cons (base +++ "_x")
                       (Cons (base +++ "_ysign")
                          (Cons (base +++ "_yvalue") Nil))
\end{verbatim}
\end{quote}

and let us remove the {\tt Bar} instance of {\tt SplitPorts}.

In the RTL, the module will now have the following ports:

\begin{center}
  \begin{tabular}{|c||c|c|c|}
    \hline
     Output &
     Input  &
     Input  &
     Input  \\
    \hline
     {\tt mv}            &
     {\tt mv\_1\_x}      &
     {\tt mv\_1\_sign}   &
     {\tt mv\_1\_yvalue} \\
     36 bits &
      8 bits &
      1 bit  &
      7 bits \\
    \hline
     {\tt mav}            &
     {\tt mav\_1\_x}      &
     {\tt mav\_1\_sign}   &
     {\tt mav\_1\_yvalue} \\
     36 bits &
      8 bits &
      1 bit  &
      7 bits \\
    \hline
  \end{tabular}
\end{center}

Each original input port has been split into three ports.

%  ----------------------------------------------------------------

\subsubsection{Example in {\BHb} of simplified canonical splitting}

\label{sec:splitports_examples_BH_simplified}

\index{SplitPorts@\te{SplitPorts}!Simplified canonical port splitting}
\index{Simplified canonical port splitting}
\index{Simplified canonical port splitting|seealso{\te{SplitPorts}}}

{\footnotesize  For a {\BSVb} version of this, please see
  Section~\ref{sec:splitports_examples_BSV_simplified}.}

{\footnotesize (This example needs import of {\tt SplitPorts} package.)}

The declarations of {\tt Foo} and {\tt Bar} as instances of type class
{\tt Splitports} for \emph{canonical} splitting, shown in
Section~\ref{sec:splitports_examples_BH} seem to have highly
predictable structure that follows the structure of the {\tt Foo} and
{\tt Bar} type declarations. The functions in the type classes
described in Section~\ref{sec:splitports_shallow_deep} provide
boilerplate shortcuts that encapsulate this predictable structure.
The instances could have been declared more simply and briefly as
follows:

Shallow-split of {\tt Foo}:
\begin{quote}
\begin{verbatim}
instance SplitPorts Foo (ShallowPortsOf Foo) where
  splitPorts   = shallowSplitPorts
  unsplitPorts = shallowUnsplitPorts
  portNames    = shallowSplitPortNames
\end{verbatim}
\end{quote}

Shallow-split of {\tt Bar}:
\begin{quote}
\begin{verbatim}
instance SplitPorts Bar (ShallowPortsOf Bar) where
  splitPorts   = shallowSplitPorts
  unsplitPorts = shallowUnsplitPorts
  portNames    = shallowSplitPortNames xBar
\end{verbatim}
\end{quote}

Deep-split of {\tt Bar}:
\begin{quote}
\begin{verbatim}
instance SplitPorts Bar (DeepPortsOf Bar) where
  splitPorts   = deepSplitPorts
  unsplitPorts = deepUnsplitPorts
  portNames    = deepSplitPortNames
\end{verbatim}
\end{quote}

NOTE: \hfill \fbox{\begin{minipage}{0.9\textwidth}\footnotesize

  Even this simple boilerplate is highly stylized, and will become
  unnecessary in a planned future enhancement of {\BLANGS}: the ``{\tt
  deriving via}'' clause.

\end{minipage}}

%  ----------------------------------------------------------------
%  ----------------------------------------------------------------
%  ----------------------------------------------------------------

\subsubsection{Examples in {\BSVb} of splitting user-defined types}


\label{sec:splitports_examples_BSV}

{\footnotesize
For {\BHb} versions of these, please see Section~\ref{sec:splitports_examples_BH}.}

% ----------------

\paragraph{Example ({\BSVb}): Baseline (before splitting) for port splitting examples}

\label{sec:splitports_examples_BSV_baseline}

\hspace*{1em}

Consider the following struct and interface types:
\begin{quote}
\begin{verbatim}
typedef struct {
   Int #(8) x;
   Int #(8) y;
} Foo
deriving (Bits);    // Represented as Bit #(16)

typedef struct {
   Vector #(3, Bool)           v;
   Tuple2 #(Bool, UInt #(16))  w;
   Foo                         z;
} Bar
deriving (Bits);    // Represented as Bit #(36)    (3+1+16+16)

interface SplitTest;
   method Bar                mv  (Foo a);
   method ActionValue #(Bar) mav (Foo a);
endinterface
\end{verbatim}
\end{quote}

If we instantiate a module {\tt mod} with this interface, the methods
may be invoked like this:

\begin{tabbing}
\hspace*{4em} \= \hspace*{2em} \= \kill
\hspace*{4em} \>               \> {\tt Foo a1, a2; Bar b1, b2;} \\
              \>               \> {\tt b1  > mod.mv  (a1)} \\
              \>  or           \> {\tt b2 <- mod.mav (a2)}
\end{tabbing}

In the RTL, the module will have the following ports:

\begin{center}
  \begin{tabular}{|c||c|}
    \hline
     Output &
     Input  \\
    \hline
     {\tt mv}    &
     {\tt mv\_a} \\
     36 bits &
     16 bits \\
    \hline
     {\tt mav}    &
     {\tt mav\_a} \\
     36 bits &
     16 bits \\
    \hline
\end{tabular}
\end{center}

% ----------------

\paragraph{Example (\BSV): shallow-split {\tt Foo}}

\label{sec:splitports_examples_BSV_Foo_Shallow}

\hspace*{1em}

{\footnotesize (This example does not need import of {\tt SplitPorts} package.)}

Let us add the following declaration to our program, making {\tt Foo}
an instance of the {\tt SplitPorts} type class:

\begin{quote}
\begin{verbatim}
typedef Tuple2 #(Port #(Int #(8)), Port #(Int #(8)))  T2;

instance SplitPorts #(Foo, T2);
   function T2 splitPorts (Foo f);
      return tuple2 (Port (f.x), Port (f.y));
   endfunction

   function Foo unsplitPorts (T2 t2);
      match { .vx, .vy } = case (t2) matches
                            { tagged Port .vx,
                              tagged Port .vy }: tuple2 (vx, vy);
                           endcase;
      return Foo { x:vx, y:vy };
   endfunction

   function portNames (_, base);
      return Cons ((base + "_x"),
                   Cons ((base + "_y"),
                         Nil));
   endfunction
endinstance
\end{verbatim}
\end{quote}

{\tt splitPorts} takes an argument {\tt f} of type {\tt Foo}, and
returns a 2-tuple of port descriptors, one per field of {\tt Foo}.
Each component applies the constructor {\tt Port} to the corresponding
component.

{\tt unsplitPorts} is the inverse: it takes a 2-tuple of port
descriptors and recombines them into a {\tt Foo} value.

{\tt portNames} returns a list of desired names for the ports, by
appending suffixes to bhe base name.

In the RTL, the module will have the following ports:

\begin{center}
  \begin{tabular}{|c||c|c|}
    \hline
     Output &
     Input  &
     Input  \\
    \hline
     {\tt mv}       &
     {\tt mv\_a\_x} &
     {\tt mv\_a\_y} \\
     36 bits &
      8 bits &
      8 bits \\
    \hline
     {\tt mav}       &
     {\tt mav\_a\_x} &
     {\tt mav\_a\_y} \\
     36 bits &
      8 bits &
      8 bits \\
    \hline
  \end{tabular}
\end{center}

The argument ports, of type {\tt Foo} have been split.

% ----------------

\paragraph{Example (\BSV): shallow-split {\tt Foo}, shallow-split {\tt Bar}}

\label{sec:splitports_examples_BSV_Foo_Shallow_Bar_Shallow}

\hspace*{1em}

{\footnotesize (This example does not need import of {\tt SplitPorts} package.)}

Suppose we add the following declaration to our program, making {\tt
  Bar} also an instance of the {\tt SplitPorts} type class:

\begin{quote}
\begin{verbatim}
typedef Tuple6 #(Port #(Bool), Port #(Bool), Port #(Bool),
                 Port #(Bool), Port #(UInt #(16)),
                 Port #(Foo))                                 T6;

instance SplitPorts #(Bar, T6);
   function T6 splitPorts (Bar b);
      match { .w1, .w2 } = b.w;
      return tuple6 (Port (b.v[0]), Port (b.v[1]), Port (b.v[2]),
                     Port (w1), Port (w2),
                     Port (b.z));
   endfunction

   function Bar unsplitPorts (T6 t6);
      match { .v0, .v1, .v2, .w1, .w2, .z }
      = case (t6) matches
           {tagged Port .v0, tagged Port .v1, tagged Port .v2,
            tagged Port .w1, tagged Port .w2,
            tagged Port .z}                     : tuple6 (v0,v1,v2,w1,w2,z);
        endcase;
      Vector #(3, Bool) v = ?;
      v[0] = v0; v[1] = v1; v[2] = v2;
      return Bar { v:v, w:tuple2 (w1,w2), z:z };
   endfunction

   function portNames (_, base);
      return Cons ((base + "_v_0"),
                   Cons ((base + "_v_1"),
                         Cons ((base + "_v_2"),
                               Cons ((base + "_w_1"),
                                     Cons ((base + "_w_2"),
                                           Cons ((base + "_z"),
                                                 Nil))))));
   endfunction
endinstance
\end{verbatim}
\end{quote}

In the RTL, the module will now have the following ports:

\begin{center}
  \begin{tabular}{|c|c|c|c|c|c||c|c|}
    \hline
     Output &
     Output &
     Output &
     Output &
     Output &
     Output &
     Input  &
     Input  \\
    \hline
     {\tt mv\_v\_0}  &
     {\tt mv\_v\_1}  &
     {\tt mv\_v\_2}  &
     {\tt mv\_w\_1}  &
     {\tt mv\_w\_2}  &
     {\tt mv\_z}     &
     {\tt mv\_a\_x}  &
     {\tt mv\_a\_y} \\
      1 bit  &
      1 bit  &
      1 bit  &
      1 bit  &
     16 bits &
     16 bits &
      8 bits &
      8 bits \\
    \hline
     {\tt mav\_v\_0}  &
     {\tt mav\_v\_1}  &
     {\tt mav\_v\_2}  &
     {\tt mav\_w\_1}  &
     {\tt mav\_w\_2}  &
     {\tt mav\_z}     &
     {\tt mav\_a\_x}  &
     {\tt mav\_a\_y} \\
      1 bit  &
      1 bit  &
      1 bit  &
      1 bit  &
     16 bits &
     16 bits &
      8 bits &
      8 bits \\
    \hline
  \end{tabular}
\end{center}

Even though each output port, of type {\tt Bar}, contains a field
``{\tt z}'' of type {\tt Foo}, that nested field is not split.
Splitting is only applied to the top-level type ({\tt Bar}), and in
its {\tt splitPorts} function we chose not to split the nested field
{\tt z}.

% ----------------

\paragraph{Example (\BSV): shallow-split {\tt Foo}, deep-split {\tt Bar}}

\label{sec:splitports_examples_BSV_Foo_Shallow_Bar_Deep}

\hspace*{1em}

{\footnotesize (This example does not need import of {\tt SplitPorts} package.)}

Suppose we change the declaration of the {\tt Bar} instance of {\tt
SplitPorts} to the following:

\begin{quote}
\begin{verbatim}
typedef Tuple7 #(Port #(Bool), Port #(Bool), Port #(Bool),
                 Port #(Bool), Port #(UInt #(16)),
                 Port #(Int #(8)),
                 Port #(Int #(8)))                            T7;

instance SplitPorts #(Bar, T7);
   function T7 splitPorts (Bar b);
      match { .w_1, .w_2 } = b.w;
      return tuple7 (Port (b.v[0]), Port (b.v[1]), Port (b.v[2]),
                     Port (w_1), Port (w_2),
                     Port (b.z.x), Port (b.z.y));
   endfunction

   function Bar unsplitPorts (T7 t7);
      match { .v_0, .v_1, .v_2, .w_1, .w_2, .z_x, .z_y }
      = case (t7) matches
           {tagged Port .v_0, tagged Port .v_1, tagged Port .v_2,
            tagged Port .w_1, tagged Port .w_2,
            tagged Port .z_x, tagged Port .z_y}: tuple7 (v_0,v_1,v_2,
                                                         w_1,w_2,
                                                         z_x,z_y);
        endcase;
      Vector #(3, Bool) v = ?;
      v[0] = v_0; v[1] = v_1; v[2] = v_2;
      return Bar { v:v, w:tuple2 (w_1,w_2), z:Foo {x:z_x, y:z_y} };
   endfunction

   function portNames (_, base);
      return Cons ((base + "_v_0"),
                   Cons ((base + "_v_1"),
                         Cons ((base + "_v_2"),
                               Cons ((base + "_w_1"),
                                     Cons ((base + "_w_2"),
                                           Cons ((base + "_z_x"),
                                                 Cons ((base + "_z_y"),
                                                       Nil)))))));
   endfunction
endinstance
\end{verbatim}
\end{quote}

In the RTL, the module will now have the following ports:

\begin{tabular}{|c|c|c|c|c|c|c||c|c|}
    \hline
     Output &
     Output &
     Output &
     Output &
     Output &
     Output &
     Output &
     Input  &
     Input  \\
    \hline
     {\tt mv\_v\_0}  &
     {\tt mv\_v\_1}  &
     {\tt mv\_v\_2}  &
     {\tt mv\_w\_1}  &
     {\tt mv\_w\_2}  &
     {\tt mv\_z\_x}  &
     {\tt mv\_z\_y}  &
     {\tt mv\_a\_x}  &
     {\tt mv\_a\_y} \\
      1 bit  &
      1 bit  &
      1 bit  &
      1 bit  &
     16 bits &
      8 bits &
      8 bits &
      8 bits &
      8 bits \\
    \hline
     {\tt mav\_v\_0}  &
     {\tt mav\_v\_1}  &
     {\tt mav\_v\_2}  &
     {\tt mav\_w\_1}  &
     {\tt mav\_w\_2}  &
     {\tt mav\_z\_x}  &
     {\tt mav\_z\_y}  &
     {\tt mav\_a\_x}  &
     {\tt mav\_a\_y} \\
      1 bit  &
      1 bit  &
      1 bit  &
      1 bit  &
     16 bits &
      8 bits &
      8 bits &
      8 bits &
      8 bits \\
    \hline
\end{tabular}

Now note that the nested field {\tt z} of type {\tt Bar} on the output
is also split, because we specfied that in the {\tt splitPorts}
function.

%  ----------------------------------------------------------------

\paragraph{Example (\BSV): custom-split {\tt Foo}}

\label{sec:splitports_examples_BSV_Foo_Custom}

\hspace*{1em}

{\footnotesize (This example does not need import of {\tt SplitPorts} package.)}

(Compare this with shallow-split {\tt Foo} in
Section:~\ref{sec:splitports_examples_BSV_Foo_Shallow}.)

Suppose we change the declaration of the {\tt Foo} instance of {\tt
SplitPorts} to the following:

\begin{quote}
\begin{verbatim}
typedef Tuple3 #(Port #(Int #(8)), Port #(Bool), Port #(Bit #(7)))  T3;

instance SplitPorts #(Foo, T3);

   function T3 splitPorts (Foo f);
      return tuple3 (Port (f.x),
                     Port (f.y > 0),
                     Port (truncate (pack (f.y))));
   endfunction

   function Foo unsplitPorts (T3 t3);
      match { .vx, .vs, .vy } = case (t3) matches
                                   {tagged Port .x,
                                    tagged Port .s,
                                    tagged Port .y}: tuple3 (x, s, y);
                                endcase;
      let uy = unpack (zeroExtend (vy));
      return Foo {x: vx,
                  y: (vs ? uy : negate (uy))};
   endfunction

   function portNames (_, base);
      return Cons ((base + "_x"),
                   Cons ((base + "_ysign"),
                         Cons ((base + "_y"),
                               Nil)));
   endfunction

endinstance
\end{verbatim}
\end{quote}

and let us remove the {\tt Bar} instance of {\tt SplitPorts}.

In the RTL, the module will now have the following ports:

\begin{center}
  \begin{tabular}{|c||c|c|c|}
    \hline
     Output &
     Input  &
     Input  &
     Input  \\
    \hline
     {\tt mv}            &
     {\tt mv\_a\_x}      &
     {\tt mv\_a\_sign}   &
     {\tt mv\_a\_yvalue} \\
     36 bits &
      8 bits &
      1 bit  &
      7 bits \\
    \hline
     {\tt mav}            &
     {\tt mav\_a\_x}      &
     {\tt mav\_a\_sign}   &
     {\tt mav\_a\_yvalue} \\
     36 bits &
      8 bits &
      1 bit  &
      7 bits \\
    \hline
  \end{tabular}
\end{center}

Each original input port has been split into three ports.

%  ----------------------------------------------------------------

\subsubsection{Example in {\BSVb} of simplified canonical splitting}

\label{sec:splitports_examples_BSV_simplified}

{\footnotesize  For a {\BHb} version of this, please see
  Section~\ref{sec:splitports_examples_BH_simplified}.}

{\footnotesize (This example needs import of {\tt SplitPorts} package.)}

The declarations of {\tt Foo} and {\tt Bar} as instances of type class
{\tt Splitports} for \emph{canonical} splitting, shown in
Section~\ref{sec:splitports_examples_BH} seem to have highly
predictable structure that follows the structure of the {\tt Foo} and
{\tt Bar} type declarations. The functions in the type classes
described in Section~\ref{sec:splitports_shallow_deep} provide
boilerplate shortcuts that encapsulate this predictable structure.
The instances could have been declared more simply and briefly as
follows:

Shallow-split of {\tt Foo}:
\begin{quote}
\begin{verbatim}
instance SplitPorts #(Foo, ShallowPortsOf #(Foo));
  function splitPorts   = shallowSplitPorts;
  function unsplitPorts = shallowUnsplitPorts;
  function portNames    = shallowSplitPortNames;
endinstance
\end{verbatim}
\end{quote}

Shallow-split of {\tt Bar}:
\begin{quote}
\begin{verbatim}
instance SplitPorts #(Bar, ShallowPortsOf #(Bar));
  function splitPorts   = shallowSplitPorts;
  function unsplitPorts = shallowUnsplitPorts;
  function portNames    = shallowSplitPortNames xBar;
\end{verbatim}
\end{quote}

Deep-split of {\tt Bar}:
\begin{quote}
\begin{verbatim}
instance SplitPorts #(Bar, DeepPortsOf #(Bar));
  function splitPorts   = deepSplitPorts;
  function unsplitPorts = deepUnsplitPorts;
  function portNames    = deepSplitPortNames;
\end{verbatim}
\end{quote}

NOTE: \hfill \fbox{\begin{minipage}{0.9\textwidth}\footnotesize

  Even this simple boilerplate is highly stylized, and will become
  unnecessary in a planned future enhancement of {\BLANGS}: the ``{\tt
  deriving via}'' clause.

\end{minipage}}

% ----------------------------------------------------------------

\subsubsection{Types for Splitting Vectors}

\label{sec:splitports_vector_types}

{\footnotesize (This example needs import of {\tt SplitVectors} package.)}

The following types are in the {\tt SplitVector} package:

\index{SplitVector@\te{SplitVector}!\te{SplitVector} type}
\index{SplitVector@\te{SplitVector}!\te{SplitVector} constructor}

\begin{quote}
  {\BHb} syntax:
  \hspace*{2em}\begin{tabular}[t]{l}
    {\tt data SplitVector n a = SplitVector (Vector n a)}
  \end{tabular}
\end{quote}

\begin{quote}
  {\BSVb} syntax:
  \hspace*{2em}\begin{tabular}[t]{l}
    {\tt typedef union tagged \{} \\
    {\tt \hspace*{2em}  Vector \#(n, a) SplitVector;} \\
    {\tt \} SplitVector \#(n, a);}
  \end{tabular}
\end{quote}

The {\tt SplitVector} type is a ``drop-in'' replacement for the {\tt Vector} type.

To split a vector method argument or result, replace the {\tt Vector}
type with {\tt SplitVector}.

%  ----------------------------------------------------------------

\paragraph{Example (\BH): splitting vectors}

\label{sec:splitports_examples_BH_vectors}

\hspace*{1em}

{\footnotesize (For a {\BSVb} version of this, please see the next section.)}

As a baseline, let us consider ports for the {\tt Vector} type.
Consider an interface like this:

\begin{quote}
\begin{verbatim}
interface SplitTest_Vec =
   mv1 :: Vector 2 (Int 8) -> Vector 3 Bool
\end{verbatim}
\end{quote}

and a module:

\begin{quote}
\begin{verbatim}
mkTest_Vec :: Module SplitTest_Vec
mkTest_Vec =
  module
    ...
    interface
      mv1 av = ... 
\end{verbatim}
\end{quote}

It will have the following input and output ports:

\begin{center}
  \begin{tabular}{|c||c|}
    \hline
     Output &
     Input  \\
    \hline
     mv1     &
     mv1\_1   \\
     3  bits &
     16 bits \\
    \hline
  \end{tabular}
\end{center}

Suppose we change the interface declaration to the following,
replacing {\tt Vector} with {\tt SplitVector}:

\begin{quote}
\begin{verbatim}
interface SplitTest_Vec =
   mv1 :: SplitVector 2 (Int 8) -> SplitVector 3 Bool
\end{verbatim}
\end{quote}

These ports will be split automatically, and the module will now have
the following input and output ports:

\begin{center}
  \begin{tabular}{|c|c|c||c|c|}
    \hline
     Output &
     Output &
     Output &
     Input  &
     Input  \\
    \hline
     mv1\_0     &
     mv1\_1     &
     mv1\_2     &
     mv1\_1\_0  &
     mv1\_1\_1  \\
     1  bits   &
     1  bits   &
     1  bits   &
     8  bits   &
     8  bits   \\
    \hline
  \end{tabular}

  {\footnotesize (The ``{\tt \_0}'', ``{\tt \_1}'' and ``{\tt \_2}''
    suffixes are the vector indexes of the corresponding components.)}

\end{center}

The method definition in the module, in more detail, will look like this:

\begin{quote}
\begin{verbatim}
mkTest_Vec :: Module SplitTest_Vec
    interface
      mv1 (SplitVector v1) = let
                               ... v1 has type Vector 2 (Int 8) ...
                               ... compute a value v2 of type Vector 3 Bool ...
                             in
                               SplitVector v2
\end{verbatim}
\end{quote}

%  ----------------------------------------------------------------

\paragraph{Example (\BSV): splitting vectors}

\label{sec:splitports_examples_BSV_vectors}

\hspace*{1em}

{\footnotesize (For a {\BHb} version of this, please see the previous section.)}

As a baseline, let us consider ports for the {\tt Vector} type.
Consider an interface like this:

\begin{quote}
\begin{verbatim}
interface SplitTest_Vec;
   method  Vector #(3, Bool) mv1 (Vector #(2, Int #(8)) av);
endinterface
\end{verbatim}
\end{quote}

and a module:

\begin{quote}
\begin{verbatim}
module mkTest_Vec1 (SplitTest_Vec);
   method mv1 (av);
      ...
      return ...;
   endmethod
endmodule
\end{verbatim}
\end{quote}

It will have the following input and output ports:

\begin{center}
  \begin{tabular}{|c||c|}
    \hline
     Output &
     Input  \\
    \hline
     mv1     &
     mv1\_av \\
     3  bits &
     16 bits \\
    \hline
  \end{tabular}
\end{center}

Suppose we change the interface declaration to the following,
replacing {\tt Vector} with {\tt SplitVector}:

\begin{quote}
\begin{verbatim}
interface SplitTest_Vec;
   method SplitVector #(3, Bool) mv1 (SplitVector #(2, Int #(8)) av);
endinterface
\end{verbatim}
\end{quote}

These ports will be split automatically, and the module will now have
the following input and output ports:

\begin{center}
  \begin{tabular}{|c|c|c||c|c|}
    \hline
     Output &
     Output &
     Output &
     Input  &
     Input  \\
    \hline
     mv1\_0      &
     mv1\_1      &
     mv1\_2      &
     mv1\_av\_0  &
     mv1\_av\_1  \\
     1  bits   &
     1  bits   &
     1  bits   &
     8  bits   &
     8  bits   \\
    \hline
  \end{tabular}

  {\footnotesize (The ``{\tt \_0}'', ``{\tt \_1}'' and ``{\tt \_2}''
    suffixes are the vector indexes of the corresponding components.)}

\end{center}

The method definition in the module, in more detail, will look like this:

\begin{quote}
\begin{verbatim}
   method SplitVector #(3, Bool) mv1 (SplitVector #(2, Int #(8)) asv);
      Vector #(2, Int #(8)) v1 = unSplitVector (asv);
      Vector #(3, Bool)     v2 = ...
      return SplitVector (v2);
   endmethod
\end{verbatim}
\end{quote}

We use the function {\tt unSplitVector} to get {\tt v1}, an equivalent
vector value of type {\tt Vector}; we compute a new value {\tt v2} of
type {\tt Vector}, as usual; finally, we apply the constructor {\tt
SplitVector} to get a result of type {\tt SplitVector}.

% ****************************************************************
